% Options for packages loaded elsewhere
% Options for packages loaded elsewhere
\PassOptionsToPackage{unicode}{hyperref}
\PassOptionsToPackage{hyphens}{url}
\PassOptionsToPackage{dvipsnames,svgnames,x11names}{xcolor}
%
\documentclass[
  12pt,
]{article}
\usepackage{xcolor}
\usepackage[margin=2.5cm]{geometry}
\usepackage{amsmath,amssymb}
\setcounter{secnumdepth}{5}
\usepackage{iftex}
\ifPDFTeX
  \usepackage[T1]{fontenc}
  \usepackage[utf8]{inputenc}
  \usepackage{textcomp} % provide euro and other symbols
\else % if luatex or xetex
  \usepackage{unicode-math} % this also loads fontspec
  \defaultfontfeatures{Scale=MatchLowercase}
  \defaultfontfeatures[\rmfamily]{Ligatures=TeX,Scale=1}
\fi
\usepackage{lmodern}
\ifPDFTeX\else
  % xetex/luatex font selection
  \setmainfont[]{Times New Roman}
  \setsansfont[]{Arial}
  \setmonofont[]{Consolas}
\fi
% Use upquote if available, for straight quotes in verbatim environments
\IfFileExists{upquote.sty}{\usepackage{upquote}}{}
\IfFileExists{microtype.sty}{% use microtype if available
  \usepackage[]{microtype}
  \UseMicrotypeSet[protrusion]{basicmath} % disable protrusion for tt fonts
}{}
\usepackage{setspace}
\makeatletter
\@ifundefined{KOMAClassName}{% if non-KOMA class
  \IfFileExists{parskip.sty}{%
    \usepackage{parskip}
  }{% else
    \setlength{\parindent}{0pt}
    \setlength{\parskip}{6pt plus 2pt minus 1pt}}
}{% if KOMA class
  \KOMAoptions{parskip=half}}
\makeatother
% Make \paragraph and \subparagraph free-standing
\makeatletter
\ifx\paragraph\undefined\else
  \let\oldparagraph\paragraph
  \renewcommand{\paragraph}{
    \@ifstar
      \xxxParagraphStar
      \xxxParagraphNoStar
  }
  \newcommand{\xxxParagraphStar}[1]{\oldparagraph*{#1}\mbox{}}
  \newcommand{\xxxParagraphNoStar}[1]{\oldparagraph{#1}\mbox{}}
\fi
\ifx\subparagraph\undefined\else
  \let\oldsubparagraph\subparagraph
  \renewcommand{\subparagraph}{
    \@ifstar
      \xxxSubParagraphStar
      \xxxSubParagraphNoStar
  }
  \newcommand{\xxxSubParagraphStar}[1]{\oldsubparagraph*{#1}\mbox{}}
  \newcommand{\xxxSubParagraphNoStar}[1]{\oldsubparagraph{#1}\mbox{}}
\fi
\makeatother


\usepackage{longtable,booktabs,array}
\usepackage{calc} % for calculating minipage widths
% Correct order of tables after \paragraph or \subparagraph
\usepackage{etoolbox}
\makeatletter
\patchcmd\longtable{\par}{\if@noskipsec\mbox{}\fi\par}{}{}
\makeatother
% Allow footnotes in longtable head/foot
\IfFileExists{footnotehyper.sty}{\usepackage{footnotehyper}}{\usepackage{footnote}}
\makesavenoteenv{longtable}
\usepackage{graphicx}
\makeatletter
\newsavebox\pandoc@box
\newcommand*\pandocbounded[1]{% scales image to fit in text height/width
  \sbox\pandoc@box{#1}%
  \Gscale@div\@tempa{\textheight}{\dimexpr\ht\pandoc@box+\dp\pandoc@box\relax}%
  \Gscale@div\@tempb{\linewidth}{\wd\pandoc@box}%
  \ifdim\@tempb\p@<\@tempa\p@\let\@tempa\@tempb\fi% select the smaller of both
  \ifdim\@tempa\p@<\p@\scalebox{\@tempa}{\usebox\pandoc@box}%
  \else\usebox{\pandoc@box}%
  \fi%
}
% Set default figure placement to htbp
\def\fps@figure{htbp}
\makeatother





\setlength{\emergencystretch}{3em} % prevent overfull lines

\providecommand{\tightlist}{%
  \setlength{\itemsep}{0pt}\setlength{\parskip}{0pt}}



 


\usepackage{polyglossia}
\setdefaultlanguage{russian}
\usepackage{csquotes}
\usepackage{caption}
\captionsetup[figure]{name=Рисунок}
\captionsetup[table]{name=Таблица}
\makeatletter
\@ifpackageloaded{caption}{}{\usepackage{caption}}
\AtBeginDocument{%
\ifdefined\contentsname
  \renewcommand*\contentsname{Table of contents}
\else
  \newcommand\contentsname{Table of contents}
\fi
\ifdefined\listfigurename
  \renewcommand*\listfigurename{List of Figures}
\else
  \newcommand\listfigurename{List of Figures}
\fi
\ifdefined\listtablename
  \renewcommand*\listtablename{List of Tables}
\else
  \newcommand\listtablename{List of Tables}
\fi
\ifdefined\figurename
  \renewcommand*\figurename{Рисунок}
\else
  \newcommand\figurename{Рисунок}
\fi
\ifdefined\tablename
  \renewcommand*\tablename{Таблица}
\else
  \newcommand\tablename{Таблица}
\fi
}
\@ifpackageloaded{float}{}{\usepackage{float}}
\floatstyle{ruled}
\@ifundefined{c@chapter}{\newfloat{codelisting}{h}{lop}}{\newfloat{codelisting}{h}{lop}[chapter]}
\floatname{codelisting}{Listing}
\newcommand*\listoflistings{\listof{codelisting}{List of Listings}}
\makeatother
\makeatletter
\makeatother
\makeatletter
\@ifpackageloaded{caption}{}{\usepackage{caption}}
\@ifpackageloaded{subcaption}{}{\usepackage{subcaption}}
\makeatother
\usepackage{bookmark}
\IfFileExists{xurl.sty}{\usepackage{xurl}}{} % add URL line breaks if available
\urlstyle{same}
\hypersetup{
  pdftitle={Отчёт по первому этапу реализации индивидуального проекта},
  pdfauthor={Лащиков Алексей Антонович},
  colorlinks=true,
  linkcolor={blue},
  filecolor={Maroon},
  citecolor={Blue},
  urlcolor={Blue},
  pdfcreator={LaTeX via pandoc}}


\title{Отчёт по первому этапу реализации индивидуального проекта}
\usepackage{etoolbox}
\makeatletter
\providecommand{\subtitle}[1]{% add subtitle to \maketitle
  \apptocmd{\@title}{\par {\large #1 \par}}{}{}
}
\makeatother
\subtitle{Размещение на Github pages заготовки для персонального сайта}
\author{Лащиков Алексей Антонович}
\date{}
\begin{document}
\maketitle

\renewcommand*\contentsname{Содержание}
{
\hypersetup{linkcolor=}
\setcounter{tocdepth}{3}
\tableofcontents
}

\setstretch{1.2}
\section{Цель
работы}\label{ux446ux435ux43bux44c-ux440ux430ux431ux43eux442ux44b}

Цель данного этапа --- развернуть заготовку персонального сайта на базе
генератора статических сайтов Hugo и опубликовать сайт на GitHub Pages с
автоматической сборкой через GitHub Actions.

\section{Задание}\label{ux437ux430ux434ux430ux43dux438ux435}

\begin{itemize}
\tightlist
\item
  Установить необходимое программное обеспечение.
\item
  Скачать шаблон темы сайта.
\item
  Разместить его на хостинге git.
\item
  Установить параметр для URLs сайта.
\item
  Разместить заготовку сайта на Github pages.
\end{itemize}

\section{Теоретическое
введение}\label{ux442ux435ux43eux440ux435ux442ux438ux447ux435ux441ux43aux43eux435-ux432ux432ux435ux434ux435ux43dux438ux435}

GitHub Pages --- сервис публикации статических сайтов, который позволяет
размещать веб-страницы напрямую из GitHub-репозитория. Hugo ---
генератор статических сайтов, который собирает проект в набор
HTML/CSS/JS-файлов. Для автоматизации сборки и публикации удобно
использовать GitHub Actions: при каждом push в ветку \texttt{main}
выполняется workflow, который устанавливает зависимости, собирает сайт и
деплоит его на Pages.

В данной работе используется starter-шаблон Wowchemy/HugoBlox на Hugo.
Для сборки темы применяется TailwindCSS, поэтому в CI требуется
установка Node.js и зависимостей (через pnpm), чтобы был доступен
\texttt{tailwindcss}.

\section{Выполнение лабораторной
работы}\label{ux432ux44bux43fux43eux43bux43dux435ux43dux438ux435-ux43bux430ux431ux43eux440ux430ux442ux43eux440ux43dux43eux439-ux440ux430ux431ux43eux442ux44b}

\subsection{Установка необходимого
ПО}\label{ux443ux441ux442ux430ux43dux43eux432ux43aux430-ux43dux435ux43eux431ux445ux43eux434ux438ux43cux43eux433ux43e-ux43fux43e}

Открыл командную строку и установил Git, Go и Node.js
(Рис.~\ref{fig-001}).

\begin{figure}

\centering{

\includegraphics[width=0.8\linewidth,height=\textheight,keepaspectratio]{image/1.png}

}

\caption{\label{fig-001}Установка Git, Go и Node.js}

\end{figure}%

Установил Hugo Extended (Рис.~\ref{fig-002}).

\begin{figure}

\centering{

\includegraphics[width=0.8\linewidth,height=\textheight,keepaspectratio]{image/2.png}

}

\caption{\label{fig-002}Установка Hugo Extended}

\end{figure}%

\subsection{Получение шаблона
сайта}\label{ux43fux43eux43bux443ux447ux435ux43dux438ux435-ux448ux430ux431ux43bux43eux43dux430-ux441ux430ux439ux442ux430}

Создал рабочую папку, склонировал starter-шаблон и установил pnpm
(Рис.~\ref{fig-003}).

\begin{figure}

\centering{

\includegraphics[width=0.8\linewidth,height=\textheight,keepaspectratio]{image/3.png}

}

\caption{\label{fig-003}Создание рабочей папки, клонирование репозитория
и установка pnpm}

\end{figure}%

Запустил локальный сервер Hugo (Рис.~\ref{fig-004}).

\begin{figure}

\centering{

\includegraphics[width=0.8\linewidth,height=\textheight,keepaspectratio]{image/4.png}

}

\caption{\label{fig-004}Запуск локального сервера}

\end{figure}%

Убедился, что сайт открывается в браузере по адресу
\texttt{http://localhost:1313/} (Рис.~\ref{fig-005}).

\begin{figure}

\centering{

\includegraphics[width=0.8\linewidth,height=\textheight,keepaspectratio]{image/5.png}

}

\caption{\label{fig-005}Демонстрация сайта}

\end{figure}%

\subsection{Создание репозитория
GitHub}\label{ux441ux43eux437ux434ux430ux43dux438ux435-ux440ux435ux43fux43eux437ux438ux442ux43eux440ux438ux44f-github}

Для публикации сайта создал публичный репозиторий с именем
(\texttt{alexeylashikov.github.io}) (Рис.~\ref{fig-006}).

\begin{figure}

\centering{

\includegraphics[width=0.8\linewidth,height=\textheight,keepaspectratio]{image/6.png}

}

\caption{\label{fig-006}Создание репозитория}

\end{figure}%

\subsection{Настройка
baseURL}\label{ux43dux430ux441ux442ux440ux43eux439ux43aux430-baseurl}

Открыл файл конфигурации \texttt{hugo.yaml} и задал параметр
\texttt{baseURL}, соответствующий будущему адресу GitHub Pages
(Рис.~\ref{fig-007}).

\begin{figure}

\centering{

\includegraphics[width=0.8\linewidth,height=\textheight,keepaspectratio]{image/7.png}

}

\caption{\label{fig-007}Настройка baseURL в hugo.yaml}

\end{figure}%

\subsection{Подготовка контента для стабильной
сборки}\label{ux43fux43eux434ux433ux43eux442ux43eux432ux43aux430-ux43aux43eux43dux442ux435ux43dux442ux430-ux434ux43bux44f-ux441ux442ux430ux431ux438ux43bux44cux43dux43eux439-ux441ux431ux43eux440ux43aux438}

Для стабильной сборки в CI временно отключил демонстрационный блог, так
как он содержал примеры с загрузкой удалённых изображений, что могло
приводить к нестабильности сборки. Для этого перенёс каталог
\texttt{content/blog} в отдельную папку \texttt{content-disabled/blog}
(Рис.~\ref{fig-008}).

\begin{figure}

\centering{

\includegraphics[width=0.8\linewidth,height=\textheight,keepaspectratio]{image/8.png}

}

\caption{\label{fig-008}Временное отключение демо-блога}

\end{figure}%

\subsection{Публикация проекта в
GitHub}\label{ux43fux443ux431ux43bux438ux43aux430ux446ux438ux44f-ux43fux440ux43eux435ux43aux442ux430-ux432-github}

Инициализировал git-репозиторий в проекте, создал коммит и отправил
проект в репозиторий на GitHub (Рис.~\ref{fig-009}).

\begin{figure}

\centering{

\includegraphics[width=0.8\linewidth,height=\textheight,keepaspectratio]{image/9.png}

}

\caption{\label{fig-009}Коммит и публикация проекта в GitHub}

\end{figure}%

\subsection{Настройка GitHub Pages и GitHub
Actions}\label{ux43dux430ux441ux442ux440ux43eux439ux43aux430-github-pages-ux438-github-actions}

В настройках репозитория открыл раздел Pages и выбрал источник
публикации GitHub Actions (Рис.~\ref{fig-010}).

\begin{figure}

\centering{

\includegraphics[width=0.8\linewidth,height=\textheight,keepaspectratio]{image/10.png}

}

\caption{\label{fig-010}Включение GitHub Pages с источником GitHub
Actions}

\end{figure}%

Далее создал workflow для автоматической сборки и деплоя. Файл
расположен по пути \texttt{.github/workflows/hugo.yml}. В workflow
добавлены шаги: установка Hugo Extended, установка Node.js, включение
\texttt{corepack}, установка зависимостей и сборка сайта командой
\texttt{pnpm\ exec\ hugo\ -\/-minify}, после чего артефакт публикуется
на GitHub Pages (Рис.~\ref{fig-011}).

\begin{figure}

\centering{

\includegraphics[width=0.8\linewidth,height=\textheight,keepaspectratio]{image/11.png}

}

\caption{\label{fig-011}Workflow GitHub Actions для сборки и публикации
Hugo-сайта}

\end{figure}%

После добавления workflow сделал коммит и отправил изменения в
репозиторий (Рис.~\ref{fig-012}).

\begin{figure}

\centering{

\includegraphics[width=0.8\linewidth,height=\textheight,keepaspectratio]{image/12.png}

}

\caption{\label{fig-012}Workflow GitHub Actions для сборки и публикации
Hugo-сайта}

\end{figure}%

\subsection{Проверка
результата}\label{ux43fux440ux43eux432ux435ux440ux43aux430-ux440ux435ux437ux443ux43bux44cux442ux430ux442ux430}

Открыл опубликованный сайт по адресу
\texttt{https://alexeylashikov.github.io/} и убедился, что сайт доступен
и корректно отображается (Рис.~\ref{fig-013}).

\begin{figure}

\centering{

\includegraphics[width=0.8\linewidth,height=\textheight,keepaspectratio]{image/13.png}

}

\caption{\label{fig-013}Опубликованный сайт на GitHub Pages}

\end{figure}%

\section{Выводы}\label{ux432ux44bux432ux43eux434ux44b}

В ходе реализации данного этапа была развёрнута заготовка персонального
сайта на базе генератора статических сайтов Hugo и опубликован сайт на
GitHub Pages с автоматической сборкой через GitHub Actions.




\end{document}
